% id: 203-22
% book: 베이직쎈 대수 · 12 수학적 귀납법
% section: 유형 127 수학적 귀납법; 부등식의 증명
% figure: none
% answer: 
% answer_source: 
% variant_level: 0 (원본 전사)
% 이 파일은 build.mjs 가 items.json 에서 만든다. 고칠 때는 items.json 을 고친다.
\smash{\raisebox{1.5mm}{\dmkichul{베이직쎈 대수 203쪽 22번}}}
다음은 $h>0$일 때, $n\ge 2$인 모든 자연수 $n$에 대하여 부등식 $(1+h)^n>1+nh$가 성립함을 수학적 귀납법으로 증명한 것이다.\exprbox{\begin{minipage}{0.96\linewidth}\raggedright\lineskip=4pt {\small\bfseries 증명}\par\smallskip (i) $n=2$일 때,\\ \hspace*{3em}(좌변)$=(1+h)^2=1+2h+h^2$, (우변)$=1+2h$\\ \hspace*{1.5em}이때 $h^2>0$이므로 \quad $1+2h+h^2>1+2h$\\ \hspace*{1.5em}따라서 주어진 부등식이 성립한다.\\ (ii) $n=k$ $(k\ge 2)$일 때 주어진 부등식이 성립한다고 가정하면 \quad $(1+h)^k>1+kh$\\ \hspace*{1.5em}위의 식의 양변에 $1+h$를 곱하면\\ \hspace*{3em}$(1+h)^{k+1}>(\fbox{㈎})(1+h)=1+\fbox{㈏}+kh^2$\\ \hspace*{5em}$>1+\fbox{㈏}$\\ \hspace*{1.5em}따라서 $n=k+1$일 때도 주어진 부등식이 성립한다.\\ (i), (ii)에서 $n\ge 2$인 모든 자연수 $n$에 대하여 주어진 부등식이 성립한다.\end{minipage}}위의 과정에서 ㈎, ㈏에 알맞은 것을 구하시오.
